\section{Marco Teórico}
\label{sec:marco_teorico}

\subsection{Memoria Asociativa Entrópica}
\label{sec:eam}

La Memoria Asociativa Entrópica (EAM) \cite{pineda2021entropic} es un modelo de memoria
que opera sobre funciones discretas.
Una \textbf{pista} es una función discreta $f : \{1, \ldots, n\} \to \{0, \ldots, m-1\}$,
representada como un vector $\mathbf{a} \in \{0, \ldots, m-1\}^n$.
La dimensión $n$ es el número de \emph{características} y $m$ es el número de valores posibles
por característica (el alfabeto).
La memoria se representa como una relación $\mathbf{R} \in \mathbb{N}^{n \times m}$,
donde $\mathbf{R}[i, j]$ cuenta cuántas veces la característica $i$ ha tomado el valor $j$ a
lo largo de todos los registros.

\subsubsection*{Operación de registro ($\lambda$)}

Dada una pista $\mathbf{a}$, se construye el patrón unitario:
\begin{equation}
    r(\mathbf{a})[i, j] = \begin{cases} 1 & \text{si } j = a_i \\ 0 & \text{en otro caso} \end{cases}
    \label{eq:patron-unitario}
\end{equation}
y se acumula en la memoria:
\begin{equation}
    \mathbf{R} \leftarrow \mathbf{R} + r(\mathbf{a})
    \label{eq:registro}
\end{equation}
La acumulación está acotada superiormente por $2^{16} - 1$ para evitar desbordamiento.

La idea central es esta: la memoria \emph{no almacena
patrones individuales} sino que acumula una \textbf{distribución de experiencias}.
Cada columna $\mathbf{R}[i, \cdot]$ se convierte progresivamente en una distribución de
frecuencias sobre los valores posibles de la característica $i$.
La memoria va formando abstracciones y prototipos a partir de esta superposición
estadística; no se imponen desde el exterior ni se calculan de forma separada.

\subsubsection*{Operación de reconocimiento ($\eta$)}

Se dice que la pista $\mathbf{a}$ está \emph{contenida} en $\mathbf{R}$ si, para toda
característica $i$, la celda correspondiente tiene masa positiva: $\mathbf{R}[i, a_i] > 0$.
La versión con tolerancia $\xi$ admite a lo más $\xi$ fallas:
\begin{equation}
    \eta(\mathbf{a}, \mathbf{R}, \xi) = \mathbf{1}\!\left[\,\bigl|\{i : \mathbf{R}[i, a_i] = 0\}\bigr| \leq \xi\,\right]
    \label{eq:reconocimiento}
\end{equation}
Adicionalmente, la operación $\eta$ devuelve un vector de \textbf{pesos por dimensión}:
\begin{equation}
    w_i = \mathbf{R}[i,\, a_i]
    \label{eq:pesos}
\end{equation}
Estos pesos reflejan cuán frecuentemente la característica $i$ ha tenido el valor $a_i$ en los
patrones registrados, y son la interfaz pista entre la memoria homo-asociativa y la
hetero-asociativa (véase la Sección~\ref{sec:hetero4d}).

El parámetro $\kappa$ define un umbral de reconocimiento relativo: se considera
reconocido si el peso medio $\bar{w} \geq \kappa \cdot \bar{\mu}$, donde $\bar{\mu}$ es
la media global de la memoria.
El parámetro $\iota$ controla la poda: en la \emph{relación operativa} de la memoria,
cualquier celda con peso $< \iota \cdot (\text{suma}/\text{celdas no cero})$ se
apaga temporalmente durante la recuperación, eliminando ruido estadístico de baja
frecuencia sin modificar la relación acumulada $\mathbf{R}$.

\subsubsection*{Operación de recuperación ($\beta$)}

La recuperación de un patrón a partir de una pista $\mathbf{a}$ procede columna a columna.
Para cada característica $i$, se construye una distribución de probabilidad sobre $\{0,\ldots,m-1\}$
modulada por una gaussiana de desviación $\sigma \cdot m$ centrada en el valor $a_i$
(cuando se dispone de la pista):
\begin{equation}
    p_i(j) \propto \mathbf{R}[i, j] \cdot \exp\!\left(-\frac{(j - a_i)^2}{2(\sigma m)^2}\right)
    \label{eq:distribucion-recuperacion}
\end{equation}
El valor recuperado $\hat{a}_i$ se obtiene \textbf{muestreando} esta distribución:
$\hat{a}_i \sim p_i(\cdot)$.

Esta variación no es un error del modelo; es parte de cómo recuerda:
\emph{recordar es muestrear}, no leer un valor determinista.
Diferentes ejecuciones de $\beta$ sobre la misma pista y la misma memoria pueden producir
resultados distintos.
El parámetro $\sigma$ controla la dispersión del muestreo: $\sigma = 0$ concentra la
masa en la pista exacta, valores mayores permiten explorar vecindades del espacio de
representación.

La \textbf{entropía} de la memoria se calcula como la media de las entropías de Shannon
por columna:
\begin{equation}
    H(\mathbf{R}) = \frac{1}{n} \sum_{i=1}^{n} \left(-\sum_{j=0}^{m-1} \hat{p}_{ij} \log_2 \hat{p}_{ij}\right)
    \label{eq:entropia}
\end{equation}
donde $\hat{p}_{ij} = \mathbf{R}[i,j] / \sum_l \mathbf{R}[i,l]$.
La entropía resume el estado de la memoria: una memoria con entropía alta ha
acumulado más variedad; una con entropía baja está dominada por pocos valores
o no ha sido suficientemente alimentada.

\subsection{Memoria Hetero-asociativa 4D}
\label{sec:hetero4d}

La memoria hetero-asociativa \cite{pineda2024hetero} relaciona dos dominios:
un dominio izquierdo con $n$ características y alfabeto de tamaño $m$, y un dominio
derecho con $p$ características y alfabeto de tamaño $q$.
La relación se almacena en un tensor
$\mathbf{R} \in \mathbb{N}^{n \times p \times (m+1) \times (q+1)}$,
donde los tamaños $(m+1)$ y $(q+1)$ incorporan \emph{centinelas} en la posición cero
para representar funciones parciales (características no definidas).

\subsubsection*{Registro hetero-asociativo}

Dada una pista izquierda $\mathbf{a} \in \{0,\ldots,m-1\}^n$ y una pista derecha
$\mathbf{b} \in \{0,\ldots,q-1\}^p$, el patrón unitario conjunto es:
\begin{equation}
    r(\mathbf{a}, \mathbf{b})[i, j, a_i, b_j] = 1
    \label{eq:registro-hetero}
\end{equation}
y se acumula de forma idéntica a la ecuación~\eqref{eq:registro}.

\subsubsection*{Proyección AND-ponderada}

Sea $\mathbf{a}$ una pista izquierda con pesos por dimensión $\mathbf{w}$ provenientes de
la operación $\eta$ de la memoria homo-asociativa del dominio izquierdo
(véase la ecuación~\eqref{eq:pesos}).
Se normalizan los pesos:
\begin{equation}
    w'_i = \frac{n \cdot w_i}{\displaystyle\sum_{k=1}^{n} w_k}
    \label{eq:normalizacion-pesos}
\end{equation}
de modo que $\frac{1}{n}\sum_i w'_i = 1$ (media unitaria).
La \textbf{proyección} $\Pi(\mathbf{a})$ sobre el dominio derecho se define por:
\begin{equation}
    \Pi[j, l] = \begin{cases}
        0 & \text{si } \exists\, i : \mathbf{R}[i, j, a_i, l] = 0 \\[4pt]
        \displaystyle\sum_{i=1}^{n} w'_i\, \mathbf{R}[i, j, a_i, l] & \text{en otro caso}
    \end{cases}
    \label{eq:proyeccion-AND}
\end{equation}
La condición de anulación es la \textbf{conjunción AND}: una celda del resultado se
anula si \emph{cualquiera} de las $n$ evidencias aporta cero.
Esto implementa \emph{containment} estricto: la proyección solo produce masa donde
todas las características de la pista tienen soporte simultáneo en la memoria.
Los pesos $w'_i$ ponderan la suma por la frecuencia con que cada característica de la pista
ha sido observada en la memoria homo-asociativa: características muy frecuentes en la pista
pesan más, lo cual permite que la memoria homo-asociativa module la hetero-asociativa
de forma nativa y justifica la arquitectura de cuatro memorias por agente.

La recuperación hetero-asociativa $\beta$ aplica muestreo por columna sobre $\Pi$ y,
cuando es necesario, una búsqueda local (\emph{sample-and-search}) para mejorar la
calidad del patrón recuperado.

\subsubsection*{Asimetría entre dominios}

La memoria homo-asociativa de un dominio modela la \emph{distribución marginal} de ese
dominio: genera pesos $w_i$ que reflejan cuán típica es una pista dentro de ese espacio.
La memoria hetero-asociativa, en cambio, modela la \emph{distribución conjunta} de dos
dominios y, por construcción, \textbf{no genera pesos por característica} del dominio derecho.
Esta asimetría es propia del modelo y justifica que cada agente en EAM-TMS requiera
\emph{cuatro} memorias: dos homo-asociativas (una por dominio, como
árbitros de pertenencia y fuente de pesos), una hetero-asociativa de contenido, y
una hetero-asociativa de directorio.

\subsection{Memoria Transactiva}
\label{sec:tmt}

El modelo de memoria transactiva de Wegner \cite{wegner1987transactive} describe cómo
los grupos desarrollan sistemas de memoria distribuida cuya capacidad excede la suma
de sus partes.
El mecanismo central es la separación entre \emph{contenido} y \emph{directorio}:
mientras cada miembro almacena información especializada, todos los miembros mantienen
un \emph{metadirectorio} de quién sabe qué.

Los tres procesos fundamentales son:

\begin{enumerate}
    \item \textbf{Actualización del directorio} (\emph{directory updating}):
    tras cada interacción de recuperación exitosa, \emph{todos} los miembros del grupo,
    incluso quienes no respondieron, actualizan su directorio anotando quién
    fue el especialista.
    Este proceso distribuye la información de especialización incluso a los miembros que
    no participaron activamente en la respuesta.

    \item \textbf{Asignación de información} (\emph{information allocation}):
    la información nueva se encamina automáticamente hacia el miembro cuya especialidad
    mejor corresponde.
    En la fase temprana este enrutamiento puede requerir un mediador que evalúe las
    capacidades de cada miembro; en la fase madura el propio directorio del emisor
    basta para determinarlo.

    \item \textbf{Coordinación de recuperación} (\emph{retrieval coordination}):
    cuando un miembro recibe una consulta fuera de su dominio, consulta su directorio
    y reenvía la consulta al especialista correspondiente sin intervención de ningún
    coordinador central.
    El \textbf{rechazo claro}, cuando ningún miembro reconoce el dominio de la
    consulta, es un comportamiento válido del grupo y evita que el sistema invente
    especialistas inexistentes.
\end{enumerate}

En la \textbf{fase temprana}, el grupo depende de un mediador: los miembros
no conocen aún las especialidades de los demás y las consultas deben ser difundidas a
todos.
La \textbf{fase madura} se alcanza cuando cada miembro ha acumulado suficiente evidencia
en su directorio para enrutar directamente: el mediador sale del circuito y las
interacciones se vuelven punto a punto.

La \textbf{entropía del directorio} resume el estado del sistema transactivo:
un directorio con entropía máxima $\log_2 K$ (siendo $K$ el número de agentes)
indica especialización perfectamente balanceada; una entropía baja indica colapso hacia
un solo agente o formación incompleta.
